% !PRINTED 260
\weberpage{260}

Eine Zahl von der Form $a\varepsilon$ soll eine \gesperrt{rein imaginäre}
\gesperrt{Zahl} heissen, zwei Zahlen der Form $a+b\varepsilon$, $a-b\varepsilon$ \gesperrt{con-}
\gesperrt{jugirt imaginäre Zahlen}. Die Uebertragung dieser Bezeichnungen
aus der Theorie der gewöhnlichen complexen Zahlen
auf die Zahlen des Körpers $\mathfrak{E}_{2,p}$ rechtfertigt sich durch die
Uebereinstimmung der Rechenregeln und kann zu keinem Missverständniss
führen, da in diesen Betrachtungen von den gewöhnlichen
complexen Zahlen nicht die Rede ist.

Weil immer, wenn $N$ ein quadratischer Nichtrest ist,
$N^{1/2(p-1)}=-1$ ist, so folgt aus (2):

\weq{3}{\varepsilon^p = -\varepsilon.}

Wendet man den binomischen Satz auf die $p^{\mathrm{te}}$ Potenz des
Binoms $\alpha=a+b\varepsilon$ an, und beachtet, dass alle Binomialco\"efficienten,
mit Ausnahme des ersten und des letzten, durch $p$
theilbar, also hier $=0$ sind, dass ferner nach dem \gesperrt{Fermat'schen}
Satze $a^p=a$, $b^p=b$ ist, so ergiebt sich nach (3):

\weq{4}{\alpha = a+b\varepsilon, \quad \alpha^p = a-b\varepsilon,}
und folglich durch Multiplication:

\weq{5}{\alpha^{p+1} = a^2 - Nb^2,}
woraus folgt, dass $\alpha^{p+1}$ immer eine reelle Zahl ist.

\gesperrt{Ist $\gamma$ eine primitive Wurzel des Körpers $\mathfrak{E}_{2,p}$ [\S.\,64, (8)],
so ist $\gamma^r$ dann und nur dann reell, wenn $r$ durch $p+1$
theilbar ist, und $\gamma^{p+1}$ ist eine primitive Wurzel der
Primzahl $p$ im gewöhnlichen Sinne}.

\gesperrt{Die nothwendige und hinreichende Bedingung
dafür, dass eine Zahl $\alpha$ in $\mathfrak{E}_{2,p}$ reell ist, besteht in der
Gleichung $\alpha^p=\alpha$}.

\gesperrt{Jede reelle Zahl ist als $(p+1)^{\mathrm{te}}$ Potenz einer Zahl
in $\mathfrak{E}_{2,p}$ darstellbar}.

Denn dass $\gamma^r$ reell ist, wenn $r$ durch $p+1$ theilbar ist,
folgt aus (5). Dass es aber nicht anders reell sein kann, ergiebt
sich daraus, dass nach dem Fermat'schen Satze jede reelle von
Null verschiedene Zahl der Bedingung $a^{p-1}=1$ genügt, dass
also, wenn $\gamma^r$ reell ist, $\gamma^{r(p-1)}=1$ sein muss, und es muss
daher $r(p-1)$ durch $p^2-1$ und folglich $r$ durch $p+1$ theilbar
sein. Da ferner $\gamma^{r(p+1)}$ nur dann $=1$ ist, wenn $r$ durch $p-1$
theilbar ist, so ist $\gamma^{p+1}$ primitive Wurzel von $p$.

\gesperrt{Auch der Begriff der Einheitswurzeln lässt sich auf
die Zahlen des Körpers $\mathfrak{E}_{2,p}$ übertragen}.
